\section{Primary crossings and a continued-fraction sequence}\label{inf:sec-arithmetic}

The real crossing orders and the integer orders are different objects.
We denote a real crossing by $\widehat p_m$ and use $\mathcal P$ for
the selected set of integers.

\subsection{A congruence approximation lemma}

\begin{lemma}[Congruence approximation]\label{inf:congruence}
For every positive irrational number $\alpha$ there are infinitely many
integer pairs $(P,Q)$ with
\begin{equation}\label{inf:congruence-target}
 P\equiv1\pmod2,\qquad Q\equiv1\pmod4,\qquad Q>0,
 \qquad Q|Q\alpha-P|<\frac92.
\end{equation}
They can be selected by the finite case rule in the proof from every
sufficiently late pair of consecutive continued-fraction convergents.
\end{lemma}
\begin{proof}
Write the convergents as $P_j^{\mathrm{cf}}/Q_j^{\mathrm{cf}}$ and put $v_j=(P_j^{\mathrm{cf}},Q_j^{\mathrm{cf}})$.
Use indices for which all the three preceding denominators are positive.
The determinant of consecutive vectors has absolute value one. Set
\[
 \rho_j=Q_{j-1}^{\mathrm{cf}}/Q_j^{\mathrm{cf}}\in(0,1),\qquad
 T_j^{\mathrm{cf}}=[a_{j+1};a_{j+2},\ldots]>1.
\]
The identity
$\alpha=(T_j^{\mathrm{cf}}P_j^{\mathrm{cf}}+P_{j-1}^{\mathrm{cf}})/(T_j^{\mathrm{cf}}Q_j^{\mathrm{cf}}+Q_{j-1}^{\mathrm{cf}})$ gives, for
$(P,Q)=Av_j+Bv_{j-1}$ and nonnegative integers $A,B$,
\begin{equation}\label{inf:cf-error}
 Q|Q\alpha-P|=\frac{(A+B\rho_j)|A-BT_j^{\mathrm{cf}}|}{T_j^{\mathrm{cf}}+\rho_j}.
\end{equation}
Denote this expression by $E(A,B)$ when the index is understood.
Directly from \eqref{inf:cf-error},
\begin{equation}\label{inf:cf-basic-bounds}
 E(1,0)<1,\quad E(1,1)<2,\quad E(2,1)<3,
 \quad E(1,2)<2(1+2\rho_j).
\end{equation}
For $E(2,1)$, split at $T_j^{\mathrm{cf}}=2$. On $1<T_j^{\mathrm{cf}}<2$ the quotient has
maximum at $T_j^{\mathrm{cf}}=1$ and is at most $(2+\rho_j)/(1+\rho_j)<2$.
On $T_j^{\mathrm{cf}}\ge2$ it is less than $2+\rho_j<3$.
The other three inequalities follow by bounding respectively
$1/(T_j^{\mathrm{cf}}+\rho_j)<1$, $(T_j^{\mathrm{cf}}-1)/(T_j^{\mathrm{cf}}+\rho_j)<1$, and
$(2T_j^{\mathrm{cf}}-1)/(T_j^{\mathrm{cf}}+\rho_j)<2$.

We first give a helper rule. Suppose $v_k$ has even numerator and odd
denominator, $v_{k-1}$ has odd numerator and even denominator, and
$Q_k^{\mathrm{cf}}+Q_{k-1}^{\mathrm{cf}}\equiv3\pmod4$. The vector
$S=3v_k+v_{k-1}$ has the target residues. Put $T=T_k^{\mathrm{cf}}$ and
$\rho=\rho_k$. If $T\ge2$, its error is less than four: for
$2\le T\le3$ it is at most $(3+\rho)/(2+\rho)<3/2$, and for
$T\ge3$ it is less than $3+\rho<4$.
If $1<T<2$, then $a_{k+1}=1$ and $v_{k+1}=v_k+v_{k-1}$.
The alternative $D=3v_{k+1}$ also has the target residues. Their
error products are
\[
 E_S=\frac{(3+\rho)(3-T)}{T+\rho},\qquad
 E_D=\frac{9(1+\rho)(T-1)}{T+\rho}.
\]
The inequalities $E_S\le3$ and $E_D\le3$ hold respectively when
\[
 T\ge\frac9{6+\rho},\qquad
 T\le\frac{3+4\rho}{2+3\rho}.
\]
These two ranges cover $(1,2)$, because
$(3+4\rho)(6+\rho)-9(2+3\rho)=4\rho^2\ge0$.
Choose the candidate with smaller error. This fixes the helper rule
and always gives error less than four, with positive denominator at
least $Q_{k-1}^{\mathrm{cf}}$.

Consider now $v_j,v_{j-1}$. There are two cases.

\emph{Case I: one vector has both entries odd.}
Call that vector $v_k$, where $k=j$ or $j-1$.
If $Q_k^{\mathrm{cf}}\equiv1\pmod4$, take $v_k$ itself.
Otherwise $Q_k^{\mathrm{cf}}\equiv3\pmod4$.
If $Q_{k-1}^{\mathrm{cf}}$ is even, its numerator is odd by the determinant condition.
For even $a_{k+1}$, take $3v_k$; here $T_k^{\mathrm{cf}}>2$ and its error is
$9/(T_k^{\mathrm{cf}}+\rho_k)<9/2$.
For odd $a_{k+1}$, the vector $v_{k+1}$ has even numerator and odd
denominator. Take $2v_{k+1}+v_k$. Its denominator is $1$ modulo four,
its numerator is odd, and \eqref{inf:cf-basic-bounds} at index $k+1$
gives error less than three.

If $Q_{k-1}^{\mathrm{cf}}$ is odd, its numerator is even.
When $a_k\ge2$ one has $\rho_k<1/2$; take
$v_k+2v_{k-1}$. Its denominator is $1$ modulo four and its error is
less than $2(1+2\rho_k)<4$.
When $a_k=1$, the vector $v_{k-2}=v_k-v_{k-1}$ has odd numerator
and even denominator. The pair $v_{k-1},v_{k-2}$ satisfies the helper
hypotheses, since its denominator sum is $Q_k^{\mathrm{cf}}\equiv3\pmod4$.
Apply the helper rule at index $k-1$.

\emph{Case II: neither vector has both entries odd.}
Their types modulo two must be (even, odd) and (odd, even): neither
vector is zero modulo two, and their determinant is odd.
If $Q_j^{\mathrm{cf}}+Q_{j-1}^{\mathrm{cf}}\equiv1\pmod4$, use their sum, whose error is less
than two. If that sum is $3$ modulo four and $Q_j^{\mathrm{cf}}$ is odd, apply the
helper rule directly.

It remains that $v_j$ has odd numerator and even denominator $Q_{\mathrm e}$,
$v_{j-1}$ has even numerator and odd denominator $Q_{\mathrm o}$, and
$Q_{\mathrm e}+Q_{\mathrm o}\equiv3\pmod4$.
Thus $Q_{\mathrm e}-Q_{\mathrm o}\equiv1\pmod4$.
Use $v_{j-2}=v_j-a_jv_{j-1}$ and distinguish the residue of $a_j$.
If $a_j\equiv1\pmod4$, take $v_{j-2}$.
If $a_j\equiv2\pmod4$, take $v_{j-1}+v_{j-2}$; its error is less
than two by \eqref{inf:cf-basic-bounds} at index $j-1$.
If $a_j\equiv3\pmod4$, the denominator of $v_{j-2}$ is $3$ modulo
four; take $2v_{j-1}+v_{j-2}$, whose error is less than three.
In these three choices the numerator is odd and the denominator is
$1$ modulo four. If $a_j\equiv0\pmod4$, the pair
$v_{j-1},v_{j-2}$ has the helper types and denominator sum $3$ modulo
four, so apply that rule at index $j-1$.

Every selected denominator is positive and at least $Q_{j-3}^{\mathrm{cf}}$.
These lower bounds tend to infinity. Thus infinitely many distinct
pairs have been constructed, and all their errors are strictly below
$9/2$.
\end{proof}

\subsection{The crossing expansion}

Define the fixed constants
\begin{equation}\label{inf:arithmetic-constants}
 c_{\mathrm{ar}}=\sqrt3-\frac\pi3,
 \qquad \alpha_{\mathrm{ar}}=\frac\pi{2c_{\mathrm{ar}}},
 \qquad d_{\mathrm{ar}}=\frac5{9\sqrt3c_{\mathrm{ar}}},
 \qquad \delta_0=1-\frac\pi{3\sqrt3}.
\end{equation}
Also write
\begin{equation}\label{inf:algebraic-radius}
 R_{\mathrm{alg}}(v)=2\sqrt{(v+1)(v+2)},\qquad
 A_m=\frac{\alpha_{\mathrm{ar}}(4m+1)-3}{2}.
\end{equation}
These constants are unrelated to the block parameters $\ab,\bb$.

\begin{proposition}[Primary-crossing expansion]\label{inf:primary-expansion}
For every sufficiently large integer $m$, there is a unique primary
crossing $\widehat p_m$ in a fixed $O(A_m^{-1})$ neighbourhood of $A_m$,
meaning
\[
 J_{\widehat p_m}(R_{\mathrm{alg}}(\widehat p_m))
 =J_{\widehat p_m+3}(R_{\mathrm{alg}}(\widehat p_m))=0.
\]
It satisfies
\begin{equation}\label{inf:root-expansion}
 \widehat p_m=A_m-\frac{d_{\mathrm{ar}}}{A_m}+O(A_m^{-2}),
\end{equation}
with an absolute remainder constant. No positive integer or
half-integer is a primary crossing.
\end{proposition}
\begin{proof}
At a zero $R$ of $J_v$, the three-term recurrence gives
\begin{equation}\label{inf:crossing-recurrence}
 \frac{J_{v+3}(R)}{J_{v-1}(R)}
 =1-\frac{4(v+1)(v+2)}{R^2}.
\end{equation}
The denominator is nonzero. Thus the crossing equation is precisely
$J_v(R_{\mathrm{alg}}(v))=0$.

At $R=R_{\mathrm{alg}}(v)$, direct Taylor expansion gives
\begin{align*}
 \sqrt{R^2-v^2}&=\sqrt3v+2\sqrt3-\frac2{\sqrt3v}+O(v^{-2}),\\
 v\arccos(v/R)&=\frac{\pi v}{3}+\frac{\sqrt3}{2}
                       -\frac{11}{4\sqrt3v}+O(v^{-2}).
\end{align*}
Consequently the phase of Lemma~\ref{inf:debye-first} is
\[
 \Xi(v,R)=c_{\mathrm{ar}}v+\frac{3\sqrt3}{2}-\frac\pi4
                         +\frac{\sqrt3}{4v}+O(v^{-2}),
\]
and its first sine coefficient is
$b(v,R)=7/(36\sqrt3v)+O(v^{-2})$.
Lemma~\ref{inf:debye-first} applies uniformly since $R/v\to2$.
In particular, the remainder is $O(v^{-2})$ for these varying
arguments.

The identity
$c_{\mathrm{ar}}A_m+3\sqrt3/2-\pi/4=\pi(m+1/2)$ is exact.
At $v=A_m\pm K/A_m$, for a sufficiently large fixed $K$, the leading
cosine changes sign with magnitude exceeding the first correction
and its remainder. The intermediate value theorem gives a zero.
Every such zero obeys
\[
 \Xi(v,R)=\pi(m+1/2)+\frac7{36\sqrt3v}+O(v^{-2}).
\]
Substitution of the phase expansion proves
\eqref{inf:root-expansion}, since
$\{7/(36\sqrt3)-\sqrt3/4\}/c_{\mathrm{ar}}=-d_{\mathrm{ar}}$.

For uniqueness, the zero lies on one of the analytic curves
$R=j_{v,k}$, because the zeros are simple. On the interval under
consideration, Lemma~\ref{inf:bessel-order} gives
$\partial_vj_{v,k}\le\pi/2$, whereas
$R'_{\mathrm{alg}}(v)=(2v+3)/\sqrt{(v+1)(v+2)}>2$.
Thus each difference $R_{\mathrm{alg}}(v)-j_{v,k}$ is strictly
increasing. Two different indices cannot cross the algebraic radius
in an interval of length $O(A_m^{-1})$: their roots are separated by
more than $\pi$, while both the algebraic radius and the zero curves
move by $O(A_m^{-1})$ there. This proves uniqueness for large $m$.

Finally, for integer $v$, two distinct nonnegative integer-order
Bessel functions have no common positive zero (Bourget's theorem).
This classical consequence of Siegel's theorem is stated in
\cite[\S10.21(i)]{inf:dlmf}, with the source
\cite[\S15.28]{inf:watson}. Applying it to orders $v$ and $v+3$
excludes an integer primary crossing. For positive half-integers,
Lemma~\ref{inf:half-integer-noncoincidence} gives the same conclusion.
\end{proof}

\begin{definition}[Selected integer sequence]\label{inf:integer-sequence}
Apply the case rule in Lemma~\ref{inf:congruence} to
$\alpha=\alpha_{\mathrm{ar}}$, with convergents indexed starting at
$0$, for each index $j\ge6$. In its helper choose the candidate with
smaller error. Remove repeated pairs and list the resulting integers
$(P-3)/2$ in increasing order. Their set is $\mathcal P$; its increasing
enumeration is denoted by $(p_m)$.
The integer $m$ in $(p_m)$ is an enumeration index, whereas the primary
crossing associated to a selected pair has index $(Q-1)/4$.
\end{definition}

\begin{definition}[Near-primary half-lattice inputs]\label{inf:input-families}
Fix an integer $m_0$ beyond which Proposition~\ref{inf:primary-expansion}
holds. For $C_{\mathrm{ap}}>0$ set
\begin{gather*}
 \mathcal P_{\mathrm{fix}}(C_{\mathrm{ap}})=
 \left\{p\in\tfrac12\Z:p>0,\ 
 \exists m\ge m_0:\ |p-\widehat p_m|<C_{\mathrm{ap}}/\widehat p_m\right\},\\
 \mathcal P_{\mathrm{bd}}=\mathcal P_{\mathrm{fix}}(31/10).
\end{gather*}
The continued-fraction set $\mathcal P$ remains the selected integer
subsequence of Definition~\ref{inf:integer-sequence}. The symbol
$p_0$ denotes the crossing associated to the current input. It is
unique for sufficiently large inputs in each fixed-width family,
since successive crossings have separation $2\alpha_{\mathrm{ar}}+o(1)$.
\end{definition}

\begin{theorem}[Near-integer primary crossings]\label{inf:near-integer}
The set $\mathcal P$ is unbounded. For every sufficiently large selected
pair $(P,Q)$, the integer $p=(P-3)/2$ and the primary crossing
$p_0=\widehat p_{(Q-1)/4}$ satisfy
\begin{equation}\label{inf:near-integer-bound}
 0<|p-p_0|<\frac{31}{10p_0}.
\end{equation}
Thus the assertion holds for every sufficiently large member of
$\mathcal P$, with its associated crossing. The threshold is implicit.
\end{theorem}
\begin{proof}
The number $\alpha_{\mathrm{ar}}$ is irrational. Indeed, a rational
value would imply
$\pi(1+2\alpha_{\mathrm{ar}}/3)=2\alpha_{\mathrm{ar}}\sqrt3$,
contradicting the transcendence of $\pi$.
For $Q=4m+1$ and $p=(P-3)/2$, the expansion
\eqref{inf:root-expansion} gives
\[
 p_0|p_0-p|\le\frac{p_0}{2}|Q\alpha_{\mathrm{ar}}-P|
              +d_{\mathrm{ar}}\frac{p_0}{A_m}+O(A_m^{-1}).
\]
Here $p_0/Q\to\alpha_{\mathrm{ar}}/2$ and $p_0/A_m\to1$.
The congruence error bound therefore yields
\begin{equation}\label{inf:arithmetic-limsup}
 \limsup p_0|p_0-p|
 \le\frac98\alpha_{\mathrm{ar}}+d_{\mathrm{ar}}<\frac{31}{10}.
\end{equation}
The final comparison follows from rational bounds. From
$\pi<22/7$ and $\sqrt3>19/11$ one obtains
$c_{\mathrm{ar}}>157/231$, $\alpha_{\mathrm{ar}}<363/157$, and
\[
 \frac98\alpha_{\mathrm{ar}}+d_{\mathrm{ar}}
 <\frac{220099}{71592}<\frac{31}{10},
 \qquad 31\cdot71592-10\cdot220099=18362>0.
\]
Bourget's theorem gives $p\ne p_0$. The denominators tend to infinity
and $P/Q\to\alpha_{\mathrm{ar}}>0$, so the selected integers are
unbounded. Distinct sufficiently large $Q$ cannot give the same
integer $p$: the corresponding numerators would be equal, while
$\alpha_{\mathrm{ar}}|Q-Q'|\le9/(2Q)+9/(2Q')$ would tend to zero.
This also removes any ambiguity in the associated crossing for large
members of the sequence.
\end{proof}

\begin{corollary}[Sparsity]\label{inf:sequence-sparsity}
The selected integer set satisfies $\#(\mathcal P\cap[1,X])=O(\log X)$.
\end{corollary}
\begin{proof}
The rule selects at most one pair at each continued-fraction index $j$,
and its denominator is at least $Q_{j-3}^{\mathrm{cf}}$. These denominators grow at
least as the Fibonacci sequence. If the selected integer is at most
$X$, its error bound gives $Q=O(X)$, since
$P=2p+3=\alpha_{\mathrm{ar}}Q+O(Q^{-1})$.
Hence only $O(\log X)$ indices can contribute. Removing duplicates
cannot increase the count.
\end{proof}

\subsection{The detuning parameter}

Continue the zero at a primary crossing $p_0$ as $R(p)=j_{p,k_0}$.
In this paragraph $p$ is real on the continuation interval; its
selected endpoint belongs to $\tfrac12\Z$. Put
\begin{equation}\label{inf:detuning-definition}
 \chi=\frac{J_{p+3}(R)}{J'_{p+3}(R)-(p-1)J_{p+3}(R)/R},
 \qquad \theta=4Rp\chi,\qquad \tau=|\theta|.
\end{equation}

\begin{lemma}[Uniform bounded detuning]\label{inf:detuning}
At every sufficiently large $p\in\mathcal P_{\mathrm{bd}}$,
\begin{equation}\label{inf:detuning-range}
 R=2p+3+O(p^{-1}),\qquad 0<|\theta|<20.
\end{equation}
More precisely, uniformly for $|p-p_0|\le31/(10p_0)$,
\begin{equation}\label{inf:detuning-asymptotic}
 \theta=(16\delta_0+o(1))p_0(p-p_0).
\end{equation}
The denominator in \eqref{inf:detuning-definition} is nonzero in this
interval for large $p_0$.
\end{lemma}
\begin{proof}
Repeated use of the Bessel recurrences at $J_p(R)=0$ gives the exact
identity
\begin{equation}\label{inf:trace-ratio-exact}
 \chi=-\frac R{2(p+1)}\frac{R^2-4(p+1)(p+2)}{2R^2-4(p+1)(p+2)}.
\end{equation}
The omitted common factor in numerator and denominator is
$J_{p-1}(R)$, which is nonzero. The denominator in the rational
expression is $4p^2+O(p)$ throughout the interval, so the physical
normal-trace denominator is also nonzero.
Lemma~\ref{inf:bessel-order}, $R/p\to2$, and $R\to\infty$ give
\[
 R'(p)=\frac{2\pi}{3\sqrt3}+o(1),\qquad
 R'_{\mathrm{alg}}(p)=2+o(1),
\]
uniformly on that interval. Since the two radii agree at $p_0$,
\[
 R(p)-R_{\mathrm{alg}}(p)=-(2\delta_0+o(1))(p-p_0).
\]
Substitute this in \eqref{inf:trace-ratio-exact} and multiply by
$4Rp$ to prove \eqref{inf:detuning-asymptotic}.
The radius assertion follows from the same comparison and
$R_{\mathrm{alg}}(p)=2p+3+O(p^{-1})$.
The inequality $\delta_0<2/5$ and the defining distance bound
give eventually
$\tau<16(2/5)(31/10)+o(1)=19.84+o(1)<20$.
The difference of the two radii has a strictly negative derivative,
because $R'\le\pi/2<2<R'_{\mathrm{alg}}$. It vanishes only at
$p=p_0$. Proposition~\ref{inf:primary-expansion} excludes that equality
on $\tfrac12\Z$. Hence $\chi\ne0$ and $\theta\ne0$ at every input.
\end{proof}

\subsection{Low angular degrees at near-primary inputs}

\begin{lemma}[A radial pole estimate]\label{inf:riccati-gap}
Let $p\ge64$, $R=R_{\mathrm{alg}}(p)$, and
$\nu=p+\ell-1<2p+3$. If $J_\nu(R)\ne0$, set
$M_\ell=J'_\nu(R)/J_\nu(R)-(p-1)/R$. Then every positive zero $z$
of $J_\nu$ satisfies
\begin{equation}\label{inf:riccati-distance}
 |z^2-R^2|\ge\frac{2(R-1)}{|M_\ell|+2}.
\end{equation}
\end{lemma}
\begin{proof}
For positive $\lambda$ away from poles define
$w(\lambda)=\sqrt\lambda J'_\nu(R\sqrt\lambda)/
J_\nu(R\sqrt\lambda)$. Bessel's equation implies
\[
 w'(\lambda)=-\frac R{2\lambda}
             \left(w(\lambda)^2+\lambda-\frac{\nu^2}{R^2}\right).
\]
Put $\eta=1/R$ and $A=\sqrt{1+\eta}$. On
$[1-\eta,1+\eta]$ one has
$|\lambda-\nu^2/R^2|\le1+\eta$, because $\nu^2<R^2+1$.
Consequently
\[
 \left|\frac d{d\lambda}\arctan\frac{w(\lambda)}A\right|
 \le\frac{RA}{2(1-\eta)}.
\]
At the first pole on either side of one the arctangent tends to an
endpoint $\pm\pi/2$. Its distance to that endpoint at $\lambda=1$
is at least $\arctan(A/|w(1)|)\ge A/(|w(1)|+A)$.
(The latter inequality follows by differentiating
$\arctan u-u/(1+u)$.) A pole in this interval is therefore at
$\lambda$-distance at least
$2(1-\eta)/(R(|w(1)|+A))$ from one.
Now $|w(1)|+A<|M_\ell|+2$, since $(p-1)/R<1/2$ and
$A<3/2$. Multiplication by $R^2$ proves
\eqref{inf:riccati-distance} for these poles. A pole outside the
interval gives $|z^2-R^2|\ge R$, which is at least the claimed
right-hand side. Poles beyond the first on either side are farther
away. This proves the bound for the entire radial spectrum.
\end{proof}

\begin{lemma}[Uniform recurrence remainder]\label{inf:low-recurrence}
At a real primary crossing $p\ge10^6$, put $s_0=R_{\mathrm{alg}}(p)/2$
and define
\[
 q_{-1}=1,\quad q_0=0,\quad
 q_{k+1}=\frac{p+k}{s_0}q_k-q_{k-1}.
\]
Let $\mathsf a_{-1}=1,\mathsf a_0=0$,
$\mathsf a_{k+1}=\mathsf a_k-\mathsf a_{k-1}$, and
$\mathsf b_{-1}=\mathsf b_0=0$,
$\mathsf b_{k+1}=\mathsf b_k-\mathsf b_{k-1}+(k-3/2)\mathsf a_k$.
Set $\mathsf e_k=q_k-\mathsf a_k-\mathsf b_k/p$ and
$\nrm{(u,v)}_\triangle=\max(|u|,|v|,|u-v|)$. Then
\begin{equation}\label{inf:low-remainder}
 \nrm{(\mathsf e_k,\mathsf e_{k-1})}_\triangle
 \le e^{k(k+3)/p}\frac{(k+2)^4}{4p^2},\qquad
 \mathsf b_{3h}=(-1)^h3h(h-1).
\end{equation}
\end{lemma}
\begin{proof}
The transformation $(u,v)\mapsto(u-v,u)$ preserves the norm
$\max(|u|,|v|,|u-v|)$. The sequence $\mathsf a_k$ has period six and norm one.
Write $\delta_k=(p+k)/s_0-1$. From
$p+3/2-s_0=1/[4(p+3/2+s_0)]$ one obtains
\[
 \left|\delta_k-\frac{k-3/2}p\right|
 \le\frac{3(k+2)}{2p^2},\qquad
 |\delta_k|\le\frac{2(k+2)}p.
\]
Summing the first-order recurrence in the invariant norm bounds its
vector by $k(k+3)/2$. Subtracting the first-order approximation from
the exact recurrence then bounds the next remainder vector by
$(1+|\delta_k|)$ times the current vector plus $(k+2)^3/p^2$.
Iteration and $1+u\le e^u$ give an exponential no greater than
$e^{k(k+3)/p}$; the sum of the cubics is at most $(k+2)^4/4$.
This proves the first bound. Three successive steps of the recurrence
for $\mathsf b_k$ give
$\mathsf b_{3(h+1)}=-\mathsf b_{3h}+6h(-1)^{h+1}$. Its initial value is zero,
so induction gives the second formula.
\end{proof}

\begin{proposition}[Low-degree exclusion]\label{inf:low-exclusion}
Let $p_0\ge10^6$ be a sufficiently large primary crossing and
$R_0=R_{\mathrm{alg}}(p_0)$. Apart from the degree-four state at $R_0$,
every even degree $\ell\ge0$ with
$\ell-1\le\sqrt{p_0}/4$ has its ball Dirichlet spectrum at distance
greater than $64$ from $R_0^2$.
At any real $p$ with $|p-p_0|\le31/(10p_0)$ and continued radius
$R=j_{p,k_0}$, those nonseed degrees have their spectrum outside
$[R^2-32,R^2+32]$. The seed's other radial states are outside this
window as well.
\end{proposition}
\begin{proof}
At the crossing the recurrence is exactly
$q_k=J_{p_0+k}(R_0)/J_{p_0-1}(R_0)$.
Put $k=\ell-1$, initially $k\ge7$ with $k^2\le p_0/16$.
The first-order contribution in the invariant vector norm is at most
$5/112$ and the remainder in \eqref{inf:low-remainder} is at most
$81/8192$; their sum is less than $1/16$.
If $3\nmid k$, this gives $|q_k|>15/16$, $|q_{k-1}|<17/16$,
and hence
\[
 |M_{k+1}|=
 \left|\frac{q_{k-1}}{q_k}-\frac{2p_0+k-1}{R_0}\right|<4.
\]
If $3\mid k$ and $k\ge15$, then
$|\mathsf b_k|=k(k-3)/3\ge4k^2/15$. The remainder is at most
$648k^4/(625p_0^2)$, which is less than $|\mathsf b_k|/(4p_0)$.
It follows that
\[
 |q_k|\ge\frac{k^2}{5p_0},\qquad |q_k|<\frac1{32},
 \qquad |M_{k+1}|\le\frac{10p_0}{k^2}.
\]
For the last estimate use $|q_{k-1}|<17/16$ and
$(2p_0+k-1)/R_0<2$.
Lemma~\ref{inf:riccati-gap} now bounds the inverse distance by
\[
 \frac{|M_{k+1}|+2}{2(R_0-1)}
 \le\frac{27}{8k^2}\le\frac{27}{1800}<\frac1{64}
\]
in the resonant range. In the nonresonant range it gives a smaller
bound for $p_0\ge10^6$.

The remaining resonant index is $k=9$.
Equation~\eqref{inf:low-remainder} gives
$|q_9|\ge(18-3661/p_0)/p_0>17.99/p_0$ and
$|q_8|\le1+44/p_0+2501/p_0^2<1.001$.
Thus $|M_{10}|<p_0/17+2$, and its inverse distance is less than
$(p_0/17+4)/(4p_0-2)<1/64$.
For $\ell=2$ use $q_1=-1,q_0=0$, giving $|M_2|<1$.
For $\ell=6$ the exact values
$q_4=(p_0+1)/s_0$ and $q_5=(p_0+4)/(p_0+2)$ give $|M_6|<2$.
For $\ell=0$ one has $M_0=0$ by \eqref{inf:ball-traces}.
These exhaust the low even degrees except $\ell=4$.
Its remaining radial states are separated by Lemma~\ref{inf:bessel-spacing}.

To pass to $p$, continue each radial zero in its order.
Both relevant orders exceed eight along the interval, so
\eqref{inf:order-lipschitz} bounds the difference of any two such
order derivatives by less than $3/5$.
Consequently $j_{p+\ell-1,k}-R(p)$ changes by less than
$93/(50p_0)<2/p_0$.
If its squared-energy offset at $p$ had absolute value at
most $32$, then $R(p)>2p_0$ and
$j_{p+\ell-1,k}+R(p)\ge(39/10)p_0$ for large $p_0$.
Its radius offset would have absolute value below $9/p_0$,
and at $p_0$ below $11/p_0$. Its squared offset at the crossing
would then be less than
\[
 \frac{11}{p_0}\left(4p_0+6+\frac{11}{p_0}\right)<55,
\]
contradicting the bound $64$. The other radial seed states remain
separated by the spacing and the same Lipschitz estimate.
\end{proof}
